% Options for packages loaded elsewhere
\PassOptionsToPackage{unicode}{hyperref}
\PassOptionsToPackage{hyphens}{url}
\documentclass[
]{article}
\usepackage{xcolor}
\usepackage[margin=1in]{geometry}
\usepackage{amsmath,amssymb}
\setcounter{secnumdepth}{-\maxdimen} % remove section numbering
\usepackage{iftex}
\ifPDFTeX
  \usepackage[T1]{fontenc}
  \usepackage[utf8]{inputenc}
  \usepackage{textcomp} % provide euro and other symbols
\else % if luatex or xetex
  \usepackage{unicode-math} % this also loads fontspec
  \defaultfontfeatures{Scale=MatchLowercase}
  \defaultfontfeatures[\rmfamily]{Ligatures=TeX,Scale=1}
\fi
\usepackage{lmodern}
\ifPDFTeX\else
  % xetex/luatex font selection
\fi
% Use upquote if available, for straight quotes in verbatim environments
\IfFileExists{upquote.sty}{\usepackage{upquote}}{}
\IfFileExists{microtype.sty}{% use microtype if available
  \usepackage[]{microtype}
  \UseMicrotypeSet[protrusion]{basicmath} % disable protrusion for tt fonts
}{}
\makeatletter
\@ifundefined{KOMAClassName}{% if non-KOMA class
  \IfFileExists{parskip.sty}{%
    \usepackage{parskip}
  }{% else
    \setlength{\parindent}{0pt}
    \setlength{\parskip}{6pt plus 2pt minus 1pt}}
}{% if KOMA class
  \KOMAoptions{parskip=half}}
\makeatother
\usepackage{color}
\usepackage{fancyvrb}
\newcommand{\VerbBar}{|}
\newcommand{\VERB}{\Verb[commandchars=\\\{\}]}
\DefineVerbatimEnvironment{Highlighting}{Verbatim}{commandchars=\\\{\}}
% Add ',fontsize=\small' for more characters per line
\usepackage{framed}
\definecolor{shadecolor}{RGB}{248,248,248}
\newenvironment{Shaded}{\begin{snugshade}}{\end{snugshade}}
\newcommand{\AlertTok}[1]{\textcolor[rgb]{0.94,0.16,0.16}{#1}}
\newcommand{\AnnotationTok}[1]{\textcolor[rgb]{0.56,0.35,0.01}{\textbf{\textit{#1}}}}
\newcommand{\AttributeTok}[1]{\textcolor[rgb]{0.13,0.29,0.53}{#1}}
\newcommand{\BaseNTok}[1]{\textcolor[rgb]{0.00,0.00,0.81}{#1}}
\newcommand{\BuiltInTok}[1]{#1}
\newcommand{\CharTok}[1]{\textcolor[rgb]{0.31,0.60,0.02}{#1}}
\newcommand{\CommentTok}[1]{\textcolor[rgb]{0.56,0.35,0.01}{\textit{#1}}}
\newcommand{\CommentVarTok}[1]{\textcolor[rgb]{0.56,0.35,0.01}{\textbf{\textit{#1}}}}
\newcommand{\ConstantTok}[1]{\textcolor[rgb]{0.56,0.35,0.01}{#1}}
\newcommand{\ControlFlowTok}[1]{\textcolor[rgb]{0.13,0.29,0.53}{\textbf{#1}}}
\newcommand{\DataTypeTok}[1]{\textcolor[rgb]{0.13,0.29,0.53}{#1}}
\newcommand{\DecValTok}[1]{\textcolor[rgb]{0.00,0.00,0.81}{#1}}
\newcommand{\DocumentationTok}[1]{\textcolor[rgb]{0.56,0.35,0.01}{\textbf{\textit{#1}}}}
\newcommand{\ErrorTok}[1]{\textcolor[rgb]{0.64,0.00,0.00}{\textbf{#1}}}
\newcommand{\ExtensionTok}[1]{#1}
\newcommand{\FloatTok}[1]{\textcolor[rgb]{0.00,0.00,0.81}{#1}}
\newcommand{\FunctionTok}[1]{\textcolor[rgb]{0.13,0.29,0.53}{\textbf{#1}}}
\newcommand{\ImportTok}[1]{#1}
\newcommand{\InformationTok}[1]{\textcolor[rgb]{0.56,0.35,0.01}{\textbf{\textit{#1}}}}
\newcommand{\KeywordTok}[1]{\textcolor[rgb]{0.13,0.29,0.53}{\textbf{#1}}}
\newcommand{\NormalTok}[1]{#1}
\newcommand{\OperatorTok}[1]{\textcolor[rgb]{0.81,0.36,0.00}{\textbf{#1}}}
\newcommand{\OtherTok}[1]{\textcolor[rgb]{0.56,0.35,0.01}{#1}}
\newcommand{\PreprocessorTok}[1]{\textcolor[rgb]{0.56,0.35,0.01}{\textit{#1}}}
\newcommand{\RegionMarkerTok}[1]{#1}
\newcommand{\SpecialCharTok}[1]{\textcolor[rgb]{0.81,0.36,0.00}{\textbf{#1}}}
\newcommand{\SpecialStringTok}[1]{\textcolor[rgb]{0.31,0.60,0.02}{#1}}
\newcommand{\StringTok}[1]{\textcolor[rgb]{0.31,0.60,0.02}{#1}}
\newcommand{\VariableTok}[1]{\textcolor[rgb]{0.00,0.00,0.00}{#1}}
\newcommand{\VerbatimStringTok}[1]{\textcolor[rgb]{0.31,0.60,0.02}{#1}}
\newcommand{\WarningTok}[1]{\textcolor[rgb]{0.56,0.35,0.01}{\textbf{\textit{#1}}}}
\usepackage{graphicx}
\makeatletter
\newsavebox\pandoc@box
\newcommand*\pandocbounded[1]{% scales image to fit in text height/width
  \sbox\pandoc@box{#1}%
  \Gscale@div\@tempa{\textheight}{\dimexpr\ht\pandoc@box+\dp\pandoc@box\relax}%
  \Gscale@div\@tempb{\linewidth}{\wd\pandoc@box}%
  \ifdim\@tempb\p@<\@tempa\p@\let\@tempa\@tempb\fi% select the smaller of both
  \ifdim\@tempa\p@<\p@\scalebox{\@tempa}{\usebox\pandoc@box}%
  \else\usebox{\pandoc@box}%
  \fi%
}
% Set default figure placement to htbp
\def\fps@figure{htbp}
\makeatother
\setlength{\emergencystretch}{3em} % prevent overfull lines
\providecommand{\tightlist}{%
  \setlength{\itemsep}{0pt}\setlength{\parskip}{0pt}}
\usepackage{bookmark}
\IfFileExists{xurl.sty}{\usepackage{xurl}}{} % add URL line breaks if available
\urlstyle{same}
\hypersetup{
  pdftitle={Informative Priors},
  hidelinks,
  pdfcreator={LaTeX via pandoc}}

\title{Informative Priors}
\author{}
\date{\vspace{-2.5em}2026-04-17}

\begin{document}
\maketitle

\subsection{Introduction}\label{introduction}

An informative prior encodes genuine pre-data knowledge: a quantitative
summary of previous studies, a mechanistic constraint that any plausible
parameter must respect, or a structured elicitation from domain experts.
Done well, informative priors stabilise inference in small samples,
integrate evidence across studies, and let a single analysis quantify
how new data update what was already known. Done badly, they paper over
weak data and make conclusions look more certain than they should. The
discipline is to be explicit about the source, justify the choice
quantitatively, and run a sensitivity analysis whenever the prior is
doing meaningful work.

\subsection{Prerequisites}\label{prerequisites}

A working understanding of Bayesian inference, the role of weakly
informative priors as a default, and basic familiarity with
meta-analytic summaries on the relevant effect scale.

\subsection{Theory}\label{theory}

Three classical sources of informative priors:

\begin{itemize}
\tightlist
\item
  \textbf{Meta-analytic priors} condense effect estimates from previous
  trials into a Normal (or log-Normal, or \(t\)) prior on the relevant
  scale: \(\beta \sim \mathcal N(\bar\beta, \tau^2)\) where
  \(\bar\beta\) and \(\tau^2\) come from a published random-effects
  meta-analysis.
\item
  \textbf{Mechanistic priors} are derived from physical, biological, or
  chemical constraints; for example, drug half-lives are positive and
  bounded above by metabolic limits, so a log-Normal or truncated prior
  is appropriate.
\item
  \textbf{Expert-elicited priors} capture structured beliefs from
  clinicians or subject-matter specialists, typically via the SHELF
  protocol or a similar formal procedure to reduce the influence of a
  single dominant expert.
\end{itemize}

A power prior or commensurate prior down-weights the prior contribution
when it conflicts with current data --- a principled way to guard
against unverified historical assumptions.

\subsection{Assumptions}\label{assumptions}

The prior information is relevant to the current analysis (same
population, same outcome definition, same scale) and is not derived from
the same data being analysed (no double-counting). Mechanistic
constraints must be physically defensible, not stylistic.

\subsection{R Implementation}\label{r-implementation}

\begin{Shaded}
\begin{Highlighting}[]
\FunctionTok{library}\NormalTok{(brms)}

\CommentTok{\# Meta{-}analytic prior: previous studies gave OR of 0.8 (95\% CI 0.7 to 0.9)}
\CommentTok{\# On log{-}OR scale: mean {-}0.223, SD \textasciitilde{} 0.05 (narrow informative prior)}
\FunctionTok{set.seed}\NormalTok{(}\DecValTok{2026}\NormalTok{)}
\NormalTok{d }\OtherTok{\textless{}{-}} \FunctionTok{data.frame}\NormalTok{(}\AttributeTok{x =} \FunctionTok{rnorm}\NormalTok{(}\DecValTok{100}\NormalTok{), }\AttributeTok{y =} \FunctionTok{rbinom}\NormalTok{(}\DecValTok{100}\NormalTok{, }\DecValTok{1}\NormalTok{, }\FloatTok{0.3}\NormalTok{))}

\NormalTok{fit\_inf }\OtherTok{\textless{}{-}} \FunctionTok{brm}\NormalTok{(y }\SpecialCharTok{\textasciitilde{}}\NormalTok{ x, }\AttributeTok{data =}\NormalTok{ d, }\AttributeTok{family =}\NormalTok{ bernoulli,}
               \AttributeTok{prior =} \FunctionTok{prior}\NormalTok{(}\FunctionTok{normal}\NormalTok{(}\SpecialCharTok{{-}}\FloatTok{0.223}\NormalTok{, }\FloatTok{0.05}\NormalTok{), }\AttributeTok{class =} \StringTok{"b"}\NormalTok{, }\AttributeTok{coef =} \StringTok{"x"}\NormalTok{),}
               \AttributeTok{chains =} \DecValTok{4}\NormalTok{, }\AttributeTok{iter =} \DecValTok{2000}\NormalTok{, }\AttributeTok{refresh =} \DecValTok{0}\NormalTok{)}

\CommentTok{\# Compare with weakly informative}
\NormalTok{fit\_weak }\OtherTok{\textless{}{-}} \FunctionTok{brm}\NormalTok{(y }\SpecialCharTok{\textasciitilde{}}\NormalTok{ x, }\AttributeTok{data =}\NormalTok{ d, }\AttributeTok{family =}\NormalTok{ bernoulli,}
                \AttributeTok{prior =} \FunctionTok{prior}\NormalTok{(}\FunctionTok{normal}\NormalTok{(}\DecValTok{0}\NormalTok{, }\FloatTok{2.5}\NormalTok{), }\AttributeTok{class =} \StringTok{"b"}\NormalTok{),}
                \AttributeTok{chains =} \DecValTok{4}\NormalTok{, }\AttributeTok{iter =} \DecValTok{2000}\NormalTok{, }\AttributeTok{refresh =} \DecValTok{0}\NormalTok{)}

\CommentTok{\# Posterior summaries}
\FunctionTok{fixef}\NormalTok{(fit\_inf); }\FunctionTok{fixef}\NormalTok{(fit\_weak)}
\end{Highlighting}
\end{Shaded}

\subsection{Output \& Results}\label{output-results}

\texttt{fixef(fit\_inf)} returns the posterior mean and credible
interval under the meta-analytic prior; \texttt{fixef(fit\_weak)} does
the same under a weakly informative prior. The two posteriors differ in
proportion to the relative precision of the prior and the data: with
\(\mathcal N(-0.223, 0.05)\) the prior carries the weight of roughly
\(1 / 0.05^2 = 400\) ``pseudo-observations'' on the log-odds scale,
dominating the small simulated dataset.

\subsection{Interpretation}\label{interpretation}

A reporting sentence: ``Using a meta-analytic prior on the log-odds
ratio derived from three previous trials
(\(\mathcal N(-0.223, 0.05^2)\), equivalent to roughly 400
pseudo-observations), the posterior estimate was \(-0.19\) (95 \%
credible interval \(-0.28\) to \(-0.10\)); under a weakly informative
\(\mathcal N(0, 2.5)\) prior, the posterior was \(-0.04\) (95 \% CrI
\(-0.43\) to \(0.34\)).'' Reporting both posteriors lets readers see
exactly how much of the conclusion is driven by the prior.

\subsection{Practical Tips}\label{practical-tips}

\begin{itemize}
\tightlist
\item
  Document the source of every informative prior in the methods section:
  cite the meta-analysis, describe the elicitation protocol, or state
  the mechanistic constraint and its justification.
\item
  Sensitivity analyses are mandatory whenever an informative prior is
  used; report the posterior under at least one weakly informative
  alternative.
\item
  Power priors with discount factor \(a_0 \in [0, 1]\) scale the prior
  likelihood and let the analyst choose how strongly historical data
  should influence the current analysis.
\item
  Commensurate priors borrow more aggressively when historical and
  current studies are similar and back off when they conflict; they are
  particularly useful in regulatory submissions.
\item
  For elicited priors, use formal tools such as the SHELF protocol;
  ad-hoc elicitation tends to underestimate uncertainty and overstate
  consensus.
\item
  In regulatory contexts (FDA, EMA), prior choice often requires
  pre-specification and agency agreement; budget time for that
  conversation before running the analysis.
\end{itemize}

\subsection{R Packages Used}\label{r-packages-used}

\texttt{brms} and \texttt{rstanarm} for placing arbitrary priors on
regression parameters; \texttt{RBesT} for meta-analytic predictive
priors and effective sample size calculations; \texttt{SHELF} for expert
elicitation; \texttt{BayesPPD} for power-prior workflows in regulatory
settings.

\subsection{For Reviewers}\label{for-reviewers}

\textbf{What to look for in a paper using this method.}

\begin{itemize}
\tightlist
\item
  \textbf{Common misapplications.}
\item
  \textbf{Diagnostics that should be reported but often aren't.}
\item
  \textbf{Red flags in tables and figures.}
\item
  \textbf{What to verify.}
\item
  \textbf{What an adequate Methods paragraph must contain.}
\end{itemize}

\subsection{See also --- labs in this
chapter}\label{see-also-labs-in-this-chapter}

\begin{itemize}
\tightlist
\item
  \href{labs/lab-bayesian-thinking-control-vs-decision-error.Rmd}{Bayesian
  thinking; α control vs decision error}
\item
  \href{labs/lab-brms-stan-loo-hierarchical-models.Rmd}{brms / Stan,
  LOO, hierarchical models}
\end{itemize}

\end{document}
